\documentclass[11pt,a4paper]{article}

\usepackage[utf8]{inputenc}
\usepackage[T1]{fontenc}
\usepackage{lmodern}
\usepackage[british]{babel}
\usepackage{microtype}
\usepackage[margin=2.5cm]{geometry}
\usepackage{amsmath,amssymb}
\usepackage{graphicx}
\usepackage{booktabs}
\usepackage{tabularx}
\usepackage{hyperref}
\usepackage[round]{natbib}
\usepackage{authblk}
\usepackage{xcolor}
\usepackage{enumitem}
\usepackage{titlesec}

\hypersetup{colorlinks=true,linkcolor=black,citecolor=black,urlcolor=blue}

\title{The phenomenology of hallucinogenic experiences: a mixed-methods comparison of psilocybin and brugmansia\\
\large with a two-stage scene-individuation and attribute-coding framework}

\author[1]{George Fejer\thanks{Correspondence: gfejer91@gmail.com}}
\author[1]{[co-authors to be listed]}
\affil[1]{[affiliation]}

\date{\today}

\begin{document}

\maketitle

\begin{abstract}
We present a phenomenological comparison of psilocybin and brugmansia hallucinatory experiences derived from human-coded trip reports. Each of 40 trip reports (20 psilocybin, 20 brugmansia) was independently coded by two hypothesis-blind raters against a 147-item hierarchical taxonomy of perceptual, cognitive, emotional, bodily, and existential features. A post-hoc consistency pipeline then reorganised the raw codings into a scene-ID-based dataset that makes inter-rater agreement separately measurable at two tiers: (i) scene individuation --- did both raters see the same passage as a hallucinatory scene --- and (ii) attribute classification --- did they tag it with the same taxonomy items. A conservative consensus filter, retaining only scenes both raters individuated and items both raters attached, operationalises the original coding instructions' Case-2 conservative rule at the dataset level and yields a high-confidence subset for between-substance comparison.
\end{abstract}

\section{Introduction}
\label{sec:introduction}

% TODO: fill in introduction
% - the phenomenology of hallucinogenic experiences as a research problem
% - why a mixed-methods approach (neurobiology + human coding + NLP)
% - serotonergic vs anticholinergic hallucinogens: distinct neurochemistry,
%   potentially distinct phenomenology
% - why this paper focuses on psilocybin vs brugmansia specifically
% - the challenge of quantifying qualitative hallucinatory content
% - preview of the scene-ID recoding framework as the paper's methodological
%   contribution

\subsection{Background}
% TODO

\subsection{Hypotheses}
% TODO

\section{Methods}
\label{sec:methods}

This section describes (i) the source corpus of trip reports, (ii) the original two-phase human coding procedure, and (iii) the post-hoc LLM-assisted consistency pipeline that converted the raw rater spreadsheets into a scene-identifier-based dataset suitable for quantitative analysis.

\subsection{Corpus and coders}
\label{sec:corpus}

The corpus comprises 40 first-person trip reports of two hallucinogenic substances: 20 reports of psilocybin (\textit{Psilocybe} spp.) and 20 reports of brugmansia (\textit{Brugmansia} spp.). Each block of 20 reports was divided into two sub-blocks of 10 reports, and each sub-block was independently coded by two hypothesis-blind raters drawn from a pool of eight coders (Table~\ref{tab:pairs}). Raters were unaware of the study hypotheses, received identical written guidelines, and had no contact with one another during coding.

\begin{table}[h]
\centering
\begin{tabular}{lll}
\toprule
Substance \& block & Coder A & Coder B \\
\midrule
Psilocybin 01--10 & Alessandra & Brendan \\
Psilocybin 11--20 & Francesco & Susana \\
Brugmansia 01--10 & Brendan & Noah \\
Brugmansia 11--20 & Alessandra & Alessio \\
\bottomrule
\end{tabular}
\caption{Coder-pair assignments across the four blocks.}
\label{tab:pairs}
\end{table}

Trip reports were sourced from public psychonaut archives and de-identified. Each report carries a unique substance-specific identifier (e.g.\ \texttt{Psilocybe\_01}), an optional self-reported dose, and a narrative of variable length (349--5\,531 words; median $\approx$1\,500).

\subsection{Coding system}
\label{sec:taxonomy}

Raters coded each report against a taxonomy of 147 items organised hierarchically into five high-level sections and 23 mid-level categories:

\begin{itemize}[itemsep=2pt]
\item \textbf{Perceptual changes}: type of visual alteration (illusion, incrusted object, detached object, full-blown immersive), visual/auditory/tactile/olfactory/gustatory/nonsensory hallucinations (each sub-classified by object class: human, artefact, animal, plant, inorganic entity, extraordinary entity), modal status (possible/probable, possible/improbable, impossible), dynamics (object constancy vs no object constancy).
\item \textbf{Sensory changes}: sensory delight (visual, auditory, haptic, erotic), synaesthesia, enhanced mental imagery.
\item \textbf{Cognitive and emotional changes}: sense of reality (lucid, delusional), amnesia, domain-specific violations (biological and psychological projection), emotion (positive, negative).
\item \textbf{Bodily and autonomic changes}: arousal (abnormally high/low), out-of-body experience, own body distortion.
\item \textbf{Existential changes}: insight (metaphysical, religious, practical, existential), ego dissolution (social, bodily, core self), sense of surprise and unexpectedness, sense of strangeness and enigma.
\end{itemize}

A subset of the taxonomy is \emph{scene-level} (hallucination modalities, alteration type, modal status, dynamics): each of these items is attached to a specific hallucinatory episode within the trip. The remainder is \emph{trip-level} (emotion, arousal, insight, ego dissolution, etc.): these items describe the trip as a whole. The ``Judgment of reality'' row, present in the original spreadsheet, was marked deprecated in the instructions and was not coded.

\subsection{Original human coding procedure}
\label{sec:human-coding}

Raters followed the procedure specified in the project's written guidelines. In outline:

\paragraph{Step 1: scene individuation.} The atomic unit of coding was the \emph{hallucinatory scene}, defined as a coherent episode in which the narrator perceives something that is not really present (a hallucination) or misperceives a real object's properties (an illusion). Per the guidelines, a scene is not defined by sentence or object boundaries: ``a scene can change a little and still remains the same scene''. Moving from a hallucinated room to a hallucinated palace with a king on a throne, by contrast, counts as a new scene. Raters parsed each report and enumerated its scenes.

\paragraph{Step 2: attribute attachment.} For each scene, the rater attached every taxonomic item that was clearly instantiated. Object-class items were counted once per scene regardless of the number of instances (3\,000 hallucinated warriors in one scene are coded as ``Human $|$ Unknown person(s)'' with count 1). If a hallucination appeared in multiple modalities (e.g.\ a visual figure who also spoke), each modality was tagged independently.

\paragraph{Step 3: trip-level items.} Items describing the trip as a whole (emotion, arousal, insight, ego dissolution, amnesia, domain-specific violations) were attached at trip level without scene attribution.

\paragraph{Conservative rule.} The guidelines specified a conservative coding discipline: if it was not clear whether an item was instantiated in the text, raters were to treat it as not instantiated and leave the cell blank. This is described in the guidelines as ``Case 2'': if unclear, do not code; only code passages that are relatively explicit.

Each rater recorded their codings in a spreadsheet with one row per taxonomy item and, for each item, an integer count of occurrences in column E and verbatim excerpts in subsequent columns.

\subsection{Post-hoc LLM-assisted consistency pipeline}
\label{sec:pipeline}

The eight rater spreadsheets represent the same 40 trips coded twice each, but their row-by-row structure makes inter-rater disagreement difficult to decompose. A disagreement can arise from (i) the two raters individuating scenes differently, (ii) agreeing on a scene but attaching different taxonomy items to it, or (iii) simply quoting different lengths of the same passage as the exemplifying excerpt. The spreadsheet format conflates these sources of disagreement into a single per-item cell.

To disentangle them we implemented a six-stage consistency pipeline in which a large language model (Claude Opus~4.6; \citealt{anthropic2024}) was used under human direction for the one step that intrinsically requires semantic judgement --- matching each rater's excerpts to a canonical set of scenes per trip --- while all other stages were deterministic and scripted. All pipeline outputs are stored in the \texttt{1.Recoding/} directory of the project repository; the canonical source spreadsheets and trip-report \texttt{.docx} files were never modified.

\paragraph{Stage 1: taxonomy extraction.} The 147-item taxonomy was parsed from one representative rater spreadsheet (the most complete; \texttt{Psilocybe\_11-20\_Susana.xlsx}) and augmented with a hierarchical annotation (\texttt{item\_id}, \texttt{parent\_id}, \texttt{depth} $\in\{0,1,2\}$, \texttt{is\_leaf}, \texttt{has\_children}, and an \texttt{is\_scene\_level} flag). Identical taxonomies were verified across all eight rater sheets (zero structural differences). The result is \texttt{data/taxonomy.csv}.

\paragraph{Stage 2: lossless raw extraction.} Each of the eight rater spreadsheets was parsed into a long-format CSV (\texttt{data/raw\_codings.csv}, $n=2\,395$ rows) with one row per (rater, trip, item, excerpt) tuple. Trip identifiers were normalised to the canonical form \texttt{Substance\_NN} (e.g.\ stripping stray hashes, whitespace, and the four blank trip-IDs in one spreadsheet were inferred by position). This CSV is a lossless mirror of the original rater data; no information is discarded.

\paragraph{Stage 3: narrative extraction.} The four trip-report \texttt{.docx} files were parsed and split by the in-text \texttt{ID:} markers into 40 per-trip plain-text narratives, stored in \texttt{data/trips/}. Per-trip word counts and doses were extracted into a \texttt{trips\_index.csv} manifest.

\paragraph{Stage 4: adjudication-worksheet generation.} For each trip, a markdown ``adjudication worksheet'' was auto-generated that displayed (a) the full trip narrative, (b) every scene-level excerpt recorded by rater A with its approximate character-span anchor in the narrative (computed via longest-common-substring and \texttt{rapidfuzz} partial-ratio fallback), and (c) the corresponding list for rater B. Excerpts attached by the same rater to multiple taxonomy items were deduplicated to a single row to make the raters' implied scene-list visible at a glance. The 40 worksheets live in \texttt{worksheets/}.

\paragraph{Stage 5: per-trip scene adjudication.} This is the stage where semantic judgement is required and where automatic matching algorithms are unsuitable: two raters' excerpts must be compared against the underlying narrative to decide whether they refer to the same hallucinatory scene. We used Claude Opus~4.6 in a structured human-in-the-loop protocol, one trip at a time:

\begin{enumerate}[itemsep=2pt]
\item The worksheet, the source narrative, and the project's coding-instruction PDFs were provided as context.
\item The model produced a YAML adjudication file per trip listing a canonical set of scenes, each with a stable scene identifier (see naming convention below), a short description, the character-span in the narrative, the worksheet-excerpt references contributed by each rater, and the taxonomic items each rater had attached.
\item The human co-author (G.F.) read each YAML against the source narrative and the worksheet and reviewed the model's unification/splitting decisions, accepting or redirecting them. Redirections were fed back into the model to refine the convention and to retrofit already-adjudicated trips where a newly-recognised pattern applied.
\end{enumerate}

A formal naming convention encodes rater status directly in each scene's identifier, making inter-rater structure queryable from the data alone. Every canonical scene has an identifier of the form \texttt{\{trip\_id\}\_S\{NN\}[.k]\_\{STATUS\}[\_MOD]}:

\begin{itemize}[itemsep=2pt]
\item \textbf{Status suffix}: \texttt{\_AB} when both raters individuated the scene, \texttt{\_A} when only rater A individuated it, \texttt{\_B} when only rater B did.
\item \textbf{Optional structural modifier}: \texttt{\_LUMPA} and \texttt{\_LUMPB} flag a parent scene where one rater lumped what the other split; corresponding child scenes carry a dotted index (\texttt{S08.1}, \texttt{S08.2}, \ldots) and the \texttt{\_SPLITB} / \texttt{\_SPLITA} modifier.
\item A \texttt{parent\_scene\_id} column in the resulting \texttt{scenes.csv} links children back to their lump.
\end{itemize}

This convention makes three distinct inter-rater disagreement tiers separately filterable from the data:

\begin{description}[itemsep=2pt]
\item[Tier 1 (scene individuation):] Did both raters individuate this scene? Answered directly by the \texttt{\_AB}/\texttt{\_A}/\texttt{\_B} suffix.
\item[Tier 2 (attribute classification):] Given that both raters individuated a scene (\texttt{\_AB}), did they attach the same taxonomy items to it?
\item[Tier 0 (binned view, analytical only):] Did both raters see \emph{anything} in this narrative region, regardless of how they chunked it? Derived at analysis time by grouping scenes with overlapping character-spans into ``bins'' (see Stage~6). This view is computed on demand and does not modify the canonical dataset.
\end{description}

The naming convention was applied retroactively: when a new pattern was discovered during adjudication of a later trip, earlier adjudications were revisited to ensure consistent application across the full dataset. The conventions are documented in \texttt{RULES.md}, which functions as an iterative living specification.

\paragraph{Stage 6: consolidation and derived tables.} A deterministic consolidation script ingests the 40 adjudication YAMLs and produces four canonical long-format CSVs plus two derived tables:

\begin{itemize}[itemsep=2pt]
\item \texttt{trips.csv} (40 rows): one row per trip with substance, block, coder pair, dose, word count.
\item \texttt{scenes.csv} (305 rows): one row per canonical scene, with \texttt{scene\_id}, \texttt{rater\_status}, canonical span, per-rater excerpt references, and any lump/split metadata.
\item \texttt{codes.csv} (1\,546 rows): one row per (scene or trip, rater, item) attachment, joined to taxonomy hierarchy.
\item \texttt{agreement\_flags.csv} (754 rows, derived): for every shared scene, every taxonomy item that at least one rater attached is listed with \texttt{A\_coded}, \texttt{B\_coded} booleans and an \texttt{agreement} $\in \{\text{AGREE}, \text{A\_ONLY}, \text{B\_ONLY}\}$ flag.
\item \texttt{rater\_style.csv} (80 rows, derived): per-trip, per-rater counts quantifying style asymmetries (scenes individuated, scene-codes attached, modal-status uses, incrusted/detached/illusion/immersive tag uses, codes-per-scene density).
\item \texttt{analysis/scene\_bins.csv} and \texttt{analysis/bins\_summary.csv} (derived, non-canonical): the binning analytical view --- scenes with overlapping canonical spans grouped into bins, with a bin-level status (both / only\_A / only\_B) and flags for granularity-recovered agreement and chunking-split patterns. These files live in \texttt{analysis/} to signal that they are a view, not a mutation of \texttt{scenes.csv}.
\end{itemize}

\subsection{Conservative consensus filter}
\label{sec:consensus-filter}

Quantitative between-substance comparisons in this paper use a conservative consensus filter that operationalises the original guidelines' Case-2 rule at the dataset level. An observation is retained if and only if:

\begin{enumerate}[itemsep=2pt]
\item \texttt{scenes.rater\_status == \text{``both''}} (both raters individuated the scene), \textbf{and}
\item \texttt{agreement\_flags.agreement == \text{``AGREE''}} (both raters attached the item to the scene).
\end{enumerate}

The justification is direct: an item coded by only one rater is, by construction, a Case-2 passage --- at least one of the two blind raters did not find the evidence clear enough to code. The guidelines instruct such passages to be treated as not instantiated. Applying this rule at the dataset level rather than per rater removes the need to arbitrate between any two raters' idiosyncratic thresholds and retains only what converged under independent application of the instructions.

Crucially, the conservative filter is an analysis-time operation. The canonical \texttt{scenes.csv} and \texttt{codes.csv} retain every scene and every code with full provenance, and analyses at more liberal thresholds (e.g.\ treating single-rater codings as exploratory observations) can be re-run against the same dataset without any re-coding.

\subsection{Diagnostic metrics}
\label{sec:diagnostics}

Inter-rater consistency is reported at two tiers:

\begin{description}[itemsep=2pt]
\item[Scene-individuation agreement] --- the proportion of canonical scenes whose \texttt{rater\_status} is ``both''.
\item[Conditional classification agreement] --- given \texttt{rater\_status == \text{``both''}}, the proportion of \texttt{agreement\_flags} rows with \texttt{agreement == \text{AGREE}}, computed separately at taxonomy depth 0 (section), 1 (category), and 2 (leaf).
\end{description}

Rater-style asymmetries are summarised by per-coder tag-usage counts in \texttt{rater\_style.csv}, which flag cases where individual raters systematically over- or under-use particular taxonomic tags across all their trips.

\subsection{Reproducibility and open data}
\label{sec:reproducibility}

The full consistency pipeline is available in the project repository under \texttt{1.Recoding/}. All stages are reproducible from the original rater spreadsheets and trip-report files: running \texttt{scripts/01\_extract\_taxonomy.py} through \texttt{scripts/07\_binning.py} in order regenerates every output file. The 40 per-trip adjudication YAMLs and their worksheets provide a full audit trail from each rater's original excerpt to its canonical scene assignment. A \texttt{RULES.md} document specifies the naming convention and the iterative rules under which new structural patterns are added. The original rater spreadsheets remain unchanged in their source directories (\texttt{Psilocybin\_source\_coding/} and \texttt{Brugmansia\_source\_coding/}).

\section{Results}
\label{sec:results}

% TODO: fill in results after agreeing on analysis plan
% The following placeholders reflect the analyses the pipeline supports.

\subsection{Corpus and coding overview}
% Summary table: per substance, n trips, total word count, n coders, n scenes individuated, n shared.

\subsection{Inter-rater consistency}
% - Scene-individuation agreement per coder pair (Table).
% - Conditional classification agreement at taxonomy depths 0, 1, 2.
% - Key finding: 81% of disagreement is attribute-level (classification), 19% individuation.
% - Binning view: confirms adjudication already captured span-overlap cases.
% - Rater-style asymmetries: per-coder tag-usage counts.

\subsection{Phenomenological profile of psilocybin}
% Consensus-filtered scene and attribute frequencies.
% - dominant alteration type
% - dominant object classes in visual hallucinations
% - auditory / tactile / nonsensory modality frequencies
% - emotion and insight profiles
% - ego dissolution frequencies

\subsection{Phenomenological profile of brugmansia}
% same structure

\subsection{Direct comparison}
% - items significantly more frequent in one substance than the other
% - qualitatively characteristic differences:
%   - psilocybin: geometric patterns, ego dissolution, metaphysical insight, extraordinary entities
%   - brugmansia: incrusted object hallucinations of humans/acquaintances,
%     delusional belief states, amnesia
% - scene-individuation rates per substance (are brugmansia trips more
%   hallucination-dense than psilocybin trips per word, etc.)

\section{Discussion}
\label{sec:discussion}

% TODO: expand during writeup. Skeleton positions follow the audit in
% 1.Recoding/AUDIT_AND_MITIGATIONS.md.

\subsection{Phenomenological findings}
% TODO: interpret the psilocybin vs brugmansia comparison once the consensus-
% filtered results table is finalised. Expected differentiators:
% - serotonergic (psilocybin) profile: geometric patterns, ego-dissolution,
%   metaphysical insight, extraordinary-entity encounters, positive-valence
%   immersive visions
% - anticholinergic (brugmansia) profile: incrusted-object hallucinations of
%   humans and acquaintances, delusional belief states, amnesia, negative-
%   valence consequences, somatic confusion
% Situate findings in the existing literature on receptor-specific hallucinogen
% phenomenology (5-HT2A vs muscarinic-antagonist profiles).

\subsection{Methodological contribution: a scene-ID recoding framework}
We present a post-hoc framework that converts rater-level spreadsheet codings into a scene-identifier dataset with explicit status suffixes (\texttt{\_AB}, \texttt{\_A}, \texttt{\_B}) and derived agreement tables. The framework makes three tiers of inter-rater agreement separately measurable — did both raters see a hallucinatory scene in the passage, did they individuate it with the same boundaries, did they attach the same attributes — and supports both strict consensus and liberal analyses by filtering on the canonical data rather than modifying it. The convention is documented iteratively in \texttt{RULES.md} and is extensible to the other ten hallucinogens in the larger project without schema changes.

\subsection{Sources of inter-rater discrepancy and mitigation}

A systematic audit of the 40-trip dataset identifies the dominant sources of disagreement. Of the 650 total disagreement events, 19\% are individuation-level (raters disagree on whether a passage is a hallucinatory scene) and 81\% are attribute-level (raters agree a scene is present but tag it differently). This indicates that the bulk of the inter-rater problem is not failure to identify hallucinatory content but variation in how that content is labelled against the taxonomy.

Within individuation disagreement, the single largest driver is a \emph{rater-threshold mismatch}: some raters (Alessandra most notably) treat brief perceptual amplifications, embodied sensations, and insight-imagery as individuated scenes, while others require a discrete object-level hallucinatory event. This asymmetry is concentrated in one coder pair (Psilocybin 01--10, Alessandra × Brendan), where Alessandra individuated 41 scenes her partner did not. The project's own guidelines (\texttt{Guidelines.pdf \S5}) mandate a conservative Case-2 rule --- if a passage is not clearly instantiated, leave it uncoded --- and a critical reading of the deviations is that Alessandra's threshold is below the guideline's intended level. Classification-level disagreement is driven principally by ambiguity in the ``Type of visual alteration'' section (illusion, incrusted, detached, immersive), stylistic over-use or under-use of ``Modal status'' and ``Incrusted object'' tags, and a scope-layer mismatch in which the same passage is coded scene-level by one rater and trip-level by the other.

We assessed eight candidate mitigation strategies (conservative consensus filter, rater reconciliation, LLM-assisted adjudication, weighted $\kappa$, third-rater tiebreaker, scope-layer normalisation, taxonomy collapse, per-rater calibration; detailed in supplementary materials). Our recommendation is a tiered application: mechanical pre-processing (scope-layer normalisation + within-rater object-row deduplication) at consolidation time; a conservative consensus filter (\texttt{rater\_status=both AND agreement=AGREE}) as the primary analysis base; a parallel liberal analysis reporting single-rater codings as exploratory observations with provenance; and, where feasible, a targeted reconciliation pass on the Alessandra × Brendan pair following the original project guidelines' prescribed method. Our view is that the framework reported here is scaffolding, not a substitute, for the reconciliation step the original guidelines foresaw.

\subsection{Limitations}
% TODO:
% - the 40-trip sample; generalisation to other hallucinogens
% - raters were not re-engaged for reconciliation; the LLM pipeline is a proxy
% - the conservative consensus filter systematically under-represents the
%   harder-to-code phenomenology (brief perceptual changes, embodied sensations,
%   insight-imagery)
% - no weighting by report length or by rater experience
% - trip reports are self-selected voluntary disclosures from psychonaut archives;
%   the corpus is not representative of the general population of hallucinogen users

\subsection{Implications and future directions}
% TODO:
% - the coding instructions could be tightened at several specific points
%   (Type-of-visual-alteration edge cases; scope-layer rules for delusional
%   beliefs and ego-dissolution; somatic phenomena)
% - a per-rater calibration exercise on a small shared-trip set before full
%   coding begins would reduce D1 threshold-mismatch at its source
% - the scene-ID framework supports natural-language-processing follow-on work
%   (scene-level embeddings; cross-substance similarity searches)
% - extend the framework to the remaining 10 hallucinogens in the corpus


\bibliographystyle{plainnat}
\bibliography{references}

\end{document}
